\begin{scanpage}{0177}{159}
Let us note that the asymptotic expressions obtained by the method of integral equations are distinguished by their considerable complexity, and frequently require tabulation of the new special functions appearing in the expressions.

In the recently issued volume of \underline{Handbuch der Physik} [50], which is devoted to diffraction theory, the complex characteristic of plane wave scattering by a strip was studied directly, omitting the calculation of the currents. For this characteristic, a singular integral equation was formulated, the solution of which was sought in the form of an asymptotic series in reciprocal powers of $\sqrt{ka}$. The first term of the series corresponds to Equations (6.14) and (6.16). The following term takes into account the interaction of the edges, and becomes infinite with the glancing incidence of a plane wave and also for observation points lying in the strip's plane. Therefore, the simple expressions obtained in [50] do not allow one to construct the complete scattering characteristic. In [50] diffraction by a disk, a sphere, and an infinite circular cylinder was investigated, and also a review of the general methods of diffraction theory and a bibliography encompassing a large number of works (mainly German and American) were given.

The book of King and Yu [59] presented (as a rule without derivation) a series of asymptotic expressions relating to a slit and a circular hole and also to other diffraction objects. Here, however, equations from which one would be able to construct the scattering characteristics of a strip and a disk with any incidence of a plane wave also are missing.

Works on diffraction by three-dimensional bodies having edges are comparatively scarce. In the paper of Siegel et al. [41], the effective scattering surface for a finite cone with the incidence of a plane wave on it along the symmetry axis is calculated from elementary arguments. The expressions obtained here do not fully characterize the fringing field, and are suitable only for sharp cones to which we already referred in \S~17. In the papers of Keller [44], the
\end{scanpage}
